% ==========================================================
%  Common preamble for the General Relativity documents
%  (IMAPP M1, Université Clermont Auvergne)
% ==========================================================
\usepackage[utf8]{inputenc}
\usepackage[T1]{fontenc}
\usepackage{lmodern}
\usepackage[english]{babel}
\usepackage[a4paper,margin=2.4cm,headheight=14pt]{geometry}
\usepackage{amsmath,amssymb,amsthm,mathtools}
\usepackage{bm}
\usepackage{enumitem}
\usepackage{booktabs}
\usepackage{array}
\usepackage{xcolor}
\usepackage{graphicx}
\usepackage{fancyhdr}
\usepackage{tikz}
\usetikzlibrary{arrows.meta,calc,decorations.pathmorphing,decorations.markings,patterns,shapes.geometric,positioning,3d,intersections}
\usepackage{pgfplots}
\pgfplotsset{compat=1.18}
\usepackage[most]{tcolorbox}
\usepackage[colorlinks=true,linkcolor=blue!50!black,citecolor=blue!50!black,urlcolor=blue!50!black]{hyperref}
\usepackage[capitalise,noabbrev]{cleveref}
\usepackage{caption}
\captionsetup{font=small,labelfont=bf}

\allowdisplaybreaks
\setlength{\parindent}{0pt}
\setlength{\parskip}{0.55em}
\numberwithin{equation}{section}

% ---------- colours ----------
\definecolor{defcol}{RGB}{0,90,160}
\definecolor{thmcol}{RGB}{160,40,40}
\definecolor{remcol}{RGB}{90,90,90}
\definecolor{excol}{RGB}{0,120,60}
\definecolor{warncol}{RGB}{200,110,0}

% ---------- boxes ----------
\tcbset{boxbase/.style={enhanced,breakable,sharp corners=south,boxrule=0.6pt,
    left=6pt,right=6pt,top=4pt,bottom=4pt,fonttitle=\bfseries,
    before skip=8pt,after skip=8pt}}
\newtcbtheorem[number within=section]{definition}{Definition}{boxbase,
    colback=defcol!4,colframe=defcol!70!black}{def}
\newtcbtheorem[use counter from=definition]{theorem}{Theorem}{boxbase,
    colback=thmcol!4,colframe=thmcol!70!black}{thm}
\newtcbtheorem[use counter from=definition]{proposition}{Proposition}{boxbase,
    colback=thmcol!4,colframe=thmcol!70!black}{prop}
\newtcbtheorem[use counter from=definition]{example}{Example}{boxbase,
    colback=excol!4,colframe=excol!70!black}{ex}
\newtcbtheorem[use counter from=definition]{exercise}{Exercise}{boxbase,
    colback=excol!4,colframe=excol!70!black}{exc}

\newtcolorbox{keyresult}[1][]{enhanced,breakable,colback=yellow!8,colframe=orange!70!black,
    boxrule=0.8pt,left=6pt,right=6pt,top=4pt,bottom=4pt,fonttitle=\bfseries,
    title={#1},before skip=8pt,after skip=8pt}
\newtcolorbox{remarkbox}[1][Remark]{enhanced,breakable,colback=remcol!5,colframe=remcol!70,
    boxrule=0.4pt,left=6pt,right=6pt,top=3pt,bottom=3pt,fonttitle=\bfseries\small,
    title={#1},before skip=6pt,after skip=6pt}
\newtcolorbox{warningbox}[1][Warning]{enhanced,breakable,colback=warncol!6,colframe=warncol!80!black,
    boxrule=0.6pt,left=6pt,right=6pt,top=3pt,bottom=3pt,fonttitle=\bfseries\small,
    title={#1},before skip=6pt,after skip=6pt}

\theoremstyle{remark}
\newtheorem*{proofsketch}{Proof sketch}

% ---------- macros ----------
\newcommand{\dd}{\mathrm{d}}
\newcommand{\pd}{\partial}
\newcommand{\R}{\mathbb{R}}
\newcommand{\M}{\mathcal{M}}
\newcommand{\Ga}[3]{\Gamma^{#1}{}_{#2#3}}
\newcommand{\Riem}[4]{R^{#1}{}_{#2#3#4}}
\newcommand{\tr}{\operatorname{tr}}
\newcommand{\sgn}{\operatorname{sgn}}
\newcommand{\GN}{G_{\mathrm N}}
\newcommand{\rs}{r_{\mathrm s}}
\newcommand{\eq}[1]{\begin{align}#1\end{align}}
\renewcommand{\Box}{\square}
\renewcommand{\vec}[1]{\bm{#1}}
\newcommand{\order}[1]{\mathcal{O}\!\left(#1\right)}
\newcommand{\half}{\tfrac12}
